\documentclass[11pt,a4paper]{article}
\usepackage[margin=18mm,top=15mm,bottom=12mm]{geometry}
\usepackage{amsmath,bm,microtype,xcolor}
\usepackage[most]{tcolorbox}
\usepackage{enumitem}
\setlist[enumerate]{itemsep=3pt,topsep=3pt,leftmargin=1.6em}
\setlist[itemize]{itemsep=2pt,topsep=2pt,leftmargin=1.3em}
\newcommand{\code}[1]{\texttt{\small #1}}
\definecolor{hl}{RGB}{0,80,150}
\pagestyle{empty}
\begin{document}

{\large\bfseries The zero-sum sheath stabiliser --- what it actually does}\\[1pt]
{\footnotesize ported from SOLPS-ITER \code{b2stbc\_stab\_coeff\_sheath\_te/ti/ni}}
\vspace{4pt}\hrule\vspace{7pt}

\textbf{1. The problem.} With $\Phi=3T_e$ at the wall, the boundary $n_e$ and $T_e$ overshoot to
\emph{negative} in a single time step, in 2--3 cells. The step is too big for that row.

\vspace{5pt}
\textbf{2. Key fact about implicit solving.} An implicit step has two separate ingredients:
\begin{center}
\begin{tabular}{ll}
\textbf{residual} $b$ & ``how wrong am I?'' \;$\rightarrow$\; decides \textbf{WHERE} the answer is ($b=0$)\\
\textbf{Jacobian} $A$ & ``how do I move?'' \;$\rightarrow$\; decides \textbf{HOW FAST} I get there\\
\end{tabular}
\end{center}
The converged state is where $b=0$. \textcolor{hl}{\textbf{$A$ does not appear in that condition.}}
So we may stiffen $A$ as much as we like --- it cannot move the answer.

\vspace{5pt}
\textbf{3. The trick.} Add to the boundary source a term that is \emph{zero when nothing is
changing}:
\[
S^{\rm stab}=\alpha\,\Pi\,\bigl(T^{\rm old}-T^{\rm new}\bigr),
\qquad \Pi=\text{the sheath flux factor already in the term}
\]
At steady state $T^{\rm new}=T^{\rm old}$, so $S^{\rm stab}\equiv 0$: it contributes
\emph{nothing} to $b$. But its derivative $\partial S^{\rm stab}/\partial T^{\rm new}=-\alpha\Pi$
is \emph{not} zero, so it adds to $A$.

\begin{tcolorbox}[colback=blue!3,colframe=hl!60,boxrule=0.6pt,left=6pt,right=6pt,top=4pt,bottom=4pt]
\textbf{4. Worked example (all of it, in three lines).} Take a scalar boundary equation with a
source $Q$ and a sink rate $k$, solved implicitly:
\[
\frac{T^{\rm new}-T^{\rm old}}{\Delta t}=Q-k\,T^{\rm new}
\qquad\Longrightarrow\qquad
\Delta T=\frac{\overbrace{Q\Delta t-k\Delta t\,T^{\rm old}}^{\textstyle \text{numerator}}}{1+k\Delta t}
\]
Now add $S^{\rm stab}=\alpha k\,(T^{\rm old}-T^{\rm new})$:
\[
\Delta T=\frac{\;Q\Delta t-k\Delta t\,T^{\rm old}\;}{1+k\Delta t+\boldsymbol{\alpha k\Delta t}}
\]
\textbf{The numerator is identical. Only the denominator grew.}
\begin{itemize}
  \item Steady state: set $\Delta T=0$ $\Rightarrow$ $T^{*}=Q/k$ \; --- \;
        \textcolor{hl}{\textbf{the same with or without $\alpha$}}.
  \item Step size: reduced by $\dfrac{1+k\Delta t}{1+k\Delta t+\alpha k\Delta t}$ ---
        \textcolor{hl}{\textbf{the overshoot is damped}}.
\end{itemize}
\end{tcolorbox}

\vspace{4pt}
\textbf{5. In the code it is one line.} Since the physical Jacobian entry carries
$(\gamma_{\rm sheath}-1)$, adding $\alpha\Pi$ to it is the same as writing
\[
(\gamma_{\rm sheath}-1+\alpha)\ \ \text{in the \textbf{Jacobian}},
\qquad
(\gamma_{\rm sheath}-1)\ \ \text{in the \textbf{residual}}.
\]
Applied to the $n_e$, $T_i$, $T_e$ boundary rows; nothing off-diagonal, no other equation.

\vspace{2pt}
\textcolor{red!60!black}{\textbf{Careful:}} scaling \emph{both} residual and Jacobian would simply
change $\gamma_{\rm sheath}$ --- a different physics problem, a different answer. The whole point
is that \textbf{only the Jacobian} is touched.

\vspace{5pt}
\textbf{6. Analogy.} It is a \textbf{damper}, not a spring. In a mass--spring system the
equilibrium is set by the spring; adding a dashpot changes only how it settles, never where.
Here $\alpha$ is the dashpot on the boundary temperature.

\vspace{5pt}
\textbf{7. Why $\alpha$ is not a fudge factor.}
$\alpha$ is dimensionless --- a multiple of the sheath transmission coefficient already present.
SOLPS ships it, recommends $\alpha=1$--$100$, and defaults to $0$; $\alpha=0$ reproduces our
previous results bit-for-bit. It damps the \textbf{transport} rows, not the momentum condition
(SOLPS ships none for that) --- matching our failure, where $n_e$ and $T_e$ go negative
\emph{together} in the same cells.

\vspace{5pt}
\textbf{8. Result.} Drift-enabled run: previously died at $\le649$ steps
$\;\rightarrow\;$ \textbf{now past 1100 and converging} (min.\ density flat, $|v_{E,n}|$ falling,
successive drops in min.\ wall $T_e$ shrinking).

\vspace{4pt}
{\footnotesize\textcolor{black!60}{Caveat: SOLPS iterates inside a time step, so there the term
vanishes before the step is accepted and is purely a convergence aid. JOREK takes one solve per
step, so $T^{\rm old}$ is the previous \emph{time step} and genuine transients are slowed too.
The stationary state is unaffected either way.}}

\end{document}
